\section{Objetivos}

\subsection{Objetivo Geral}

Propor uma solução computacional para o problema de alocação semestral de grades horárias em Instituições de Ensino Superior por meio do IGRIS (\textit{Intelligent Grade Integration System}), um sistema modular de suporte à decisão baseado em Algoritmos Genéticos. A solução busca automatizar a geração de grades acadêmicas, reduzir conflitos de alocação, apoiar o trabalho dos coordenadores acadêmicos e permitir a avaliação comparativa de diferentes implementações evolutivas aplicadas ao \textit{University Course Timetabling Problem} (UCTP).

\subsection{Objetivos Específicos}

Os objetivos específicos do presente trabalho são:

\begin{itemize}
\item Realizar revisão bibliográfica sobre o UCTP, otimização combinatória, heurísticas, meta-heurísticas e Algoritmos Genéticos aplicados à geração de grades horárias acadêmicas.

\item Modelar uma base de dados relacional capaz de representar as principais entidades acadêmicas envolvidas no processo de alocação, como professores, disciplinas, turmas, matrizes curriculares, disponibilidades e grades horárias.

\item Definir as restrições acadêmicas e operacionais consideradas na geração das grades horárias, contemplando conflitos de professores, choques de turmas, disponibilidades docentes e critérios de qualidade da solução.

\item Desenvolver módulos de gestão de entidades acadêmicas, permitindo o cadastro, a manutenção e a consulta das informações necessárias à geração automatizada das grades horárias.

\item Implementar variações de Algoritmos Genéticos aplicadas ao UCTP, contemplando diferentes paradigmas de programação e abordagens de implementação.

\item Integrar as implementações algorítmicas ao sistema IGRIS por meio de uma arquitetura modular com adaptadores de microsserviços, permitindo a execução, substituição e comparação dos algoritmos sem comprometer os demais módulos da aplicação.

\item Desenvolver uma interface web responsiva para configuração dos dados acadêmicos, parametrização dos algoritmos e visualização das grades horárias geradas.

\item Avaliar o desempenho das implementações por meio de \textit{benchmarks} controlados, considerando métricas como latência, consumo de memória RAM, conflitos de alocação, disciplinas não alocadas e qualidade da solução gerada.
\end{itemize}
